\section{Eight Integrated Practices}
\label{sec:eight}

The contribution of F(AI)\textsuperscript{2}R is the joint adoption of
the eight practices below. Each is individually unoriginal; the
literature for any single one is older than this paper. The argument
is that adopting any one of them in isolation is insufficient against
the failure modes catalogued in Section~\ref{sec:failure-modes}, while
adopting all eight, with the pipeline of \S\ref{sec:pipeline} enforcing
the seams, is sufficient against most of them and honest about the
rest.

\subsection{Practice 1: Transcript preservation}
\label{sec:practice-transcripts}

Every working session that produces a non-trivial change to the
manuscript exports its conversation transcript and commits it under
\texttt{doc/transcripts/YYYY-MM-DD-<slug>.md}. The transcript is the
primary record of what was actually said, by which agent, in what
order. Identifiers carried with each transcript: filename, SHA-256 of
the file, commit SHA at which it was committed.

The practice attacks the most-easily-lost piece of LLM-assisted work:
the prompt-and-response sequence that produced a particular paragraph.
A reader who later asks ``which model wrote this sentence, and what
preceded it?'' has a deterministic answer.

The practice is unoriginal: open-source projects have committed
chat transcripts for years. F(AI)\textsuperscript{2}R is the
commitment to do so \emph{by default} and to bind the transcript to
the manuscript via the provenance graph: each
\texttt{prov:Activity} that authored a section carries a
\texttt{prov:used} pointing at the transcript entity.

\subsection{Practice 2: Verification-status labelling}
\label{sec:practice-ladder}

Every claim in the manuscript carries one of the five rungs defined
in \S\ref{sec:ladder}. The rung is recorded in the provenance graph,
not in the prose: a \LaTeX{} reader sees only the citation; a graph
reader sees the rung. A claim invoked at a rung lower than its evidence
deserves is allowed; a claim invoked at a rung higher than its evidence
permits is a defect surfaced by the Aligner.

The practice attacks claim--evidence drift. Without a rung, ``cited
but unread'' and ``cited and read'' look the same in the manuscript
and the difference is invisible at review time. The rung makes the
difference machine-checkable.

\subsection{Practice 3: Per-claim provenance maps}
\label{sec:practice-provenance}

Every named claim in the manuscript has a corresponding
\texttt{fair2r:Claim} IRI in \texttt{doc/provenance.ttl} with at minimum
\texttt{prov:wasGeneratedBy} an authoring activity,
\texttt{prov:wasAttributedTo} an agent (human or AI),
\texttt{prov:wasDerivedFrom} the source that supports it (if any), and
\texttt{fair2r:verificationState} one of the rungs above. The graph
is the spine of the audit pass: a third party can replay the chain
from claim to source without re-reading the prose.

The practice attacks two failure modes simultaneously: fabricated
citations (a claim with no \texttt{prov:wasDerivedFrom} cannot pass
the Aligner) and silent retraction (an invalidated claim must carry a
\texttt{prov:Invalidation} activity, not a deletion).

\subsection{Practice 4: Mirror discipline}
\label{sec:practice-mirror}

The mirror discipline of \S\ref{sec:mirror} is the fourth practice in
its own right. The pair (source, render) is the unit; the build either
agrees or fails. The condensed manuscript, the public site, and the
provenance tables are all renders that must agree with their sources.

The practice attacks placeholder dishonesty (a render that hides
unfilled material) and condensation drift (claims that survived a
length cut but lost their IRIs). \emph{Asset honesty over asset
completeness}: if an asset does not exist, ship a labelled placeholder
that visibly identifies itself as AI-authored and pending, rather than
a polished surface that conceals which load-bearing parts remain
unverified.

\subsection{Practice 5: Recursive meta-process documentation}
\label{sec:practice-recursive}

This paper is its own case study. The pipeline that produced it is the
pipeline it describes. The provenance graph for this paper is part of
this paper's release. The eight practices listed here were developed
on a different paper (the source repository
\texttt{noheton/Obscurity-Is-Dead}, abstracted here so that none of its
domain content travels with the methodology) and consolidated for
publication on this one.

The practice attacks the methodological-tract failure mode --- a paper
that prescribes a discipline it did not itself follow. Any reader can
clone the repository, rebuild the PDF, walk the graph, and verify that
each claim carrying a particular rung has the artefacts the rung
demands.

\subsection{Practice 6: Base-rate-anchored AI disclosure}
\label{sec:practice-disclosure}

Disclosure is mandatory at most venues and largely ritualised. Most
disclosure boilerplate names a tool and says it was used. The sixth
practice is to anchor disclosure to numeric base rates: when we say a
claim was \textsc{ai-confirmed}, we say so against the published rates
for the failure mode that rung is supposed to mitigate. For
fabricated citations specifically, we cite the figure that motivated
the rung's existence~\cite{walters2023fabrication,magesh2024legal} and
declare what fraction of the manuscript was \emph{not} elevated past
that rung at submission time.

The practice attacks ``perfunctory disclosure'': a sentence in the
acknowledgements that names the model and stops there. Anchored
disclosure says: this many claims sit at \textsc{ai-confirmed} and not
above; here is the literature on the residual risk; here is the
mitigation we adopted.

\subsection{Practice 7: Legal honesty about authorship}
\label{sec:practice-legal}

AI is acknowledged but is not an author. The legal position is
consistent across the jurisdictions we have surveyed: in
Germany (\S2 \emph{UrhG}), the United States (USCO 2023 guidance),
and at the major medical journals via the
\textsc{ICMJE}~\cite{icmje2023}, AI cannot hold copyright, cannot
consent to publication, and cannot be held accountable. The
\texttt{CITATION.cff} therefore lists human authors only; the
\emph{Acknowledgements} section names the AI tools used; the
provenance graph attributes activities to AI agents as
\texttt{prov:Agent} but never to a person who is also an author. The
distinction is not cosmetic. An AI that is listed as author would, on
some readings, be claimed by the human author as a copyrightable
contribution they did not in fact make.

The practice attacks two failure modes. First, the inflation of an AI
acknowledgement into a co-author entry --- the kind of move that turns
an honest disclosure into a misrepresentation. Second, the silent
treatment: an acknowledgement that does not appear at all, leaving
readers to infer the AI's contribution from style alone.

\subsection{Practice 8: FAIR alignment as precondition}
\label{sec:practice-fair}

\emph{Do not retrofit FAIR.} If the eight practices above are adopted
on day one, the four FAIR axes follow from them at low marginal cost.
If they are adopted at submission time, the cost dominates: the
manuscript has to be re-decomposed into claims, the transcripts have to
be reconstructed if they were not preserved, the verification rungs
have to be assigned in retrospect (with the verification overhead now
maximally expensive), and the prompt set has to be reverse-engineered
from session logs.

The practice is the easiest to adopt and the easiest to skip. It is
also the one that gates the others: a paper whose authors did not adopt
F(AI)\textsuperscript{2}R from day one has, in effect, given up the
audit pass on F, A, I, and R alike. We treat it as a precondition.
